%% ---------------------------------------------------------------------------
\section{The First Rows After the Reflection}\label{sec:early}

The five-state rules that come closest to minimal time fail in the first rows
of their reflected triangles: some cell of a long line fires a few steps after
the front has reached the right end (\cref{prop:nearmiss,fig:delta14}). The
condition that this never happens concerns infinitely many lengths. We show
that, for a given half-line, it reduces to finitely many lengths, and that a
certificate for the reduction can be computed by simulating the half-line.

\subsection{Finitely Many Early Rows}

For $T\ge0$ let $E(T):=\{(\tau,\kappa):0\le\kappa\le\tau\le T\}$ be the first
$T+1$ rows of a reflected triangle, that is, its cells at the relative times
$0,\dots,T$, during the first $T+1$ steps after the front reaches the right
end. If $n\ge T+2$, then $E(T)\subseteq R_n$,
and no cell of $E(T)$ lies in the firing row $\tau=n-1$.

\begin{lemma}[early rows]\label{lem:early}
Let $\delta$ be a rule, $T\ge0$, and $n,n'\ge T+2$ such that the partial values
$J_n(\kappa)$ and $J_{n'}(\kappa)$ coincide for $0\le\kappa\le T+1$. Then for
every $(\tau,\kappa)\in E(T)$ the values $R_n(\tau,\kappa)$ and
$R_{n'}(\tau,\kappa)$ are both undefined or both defined and equal.
\end{lemma}

\begin{proof}
By \cref{lem:determinacy} we may compare $\Psi(J_n,n;\tau,\kappa)$ and
$\Psi(J_{n'},n';\tau,\kappa)$. Unfolding the recursion of
\cref{lem:determinacy} from $(\tau,\kappa)\in E(T)$ visits only pairs
$(\tau'',\kappa'')$ with $0\le\kappa''\le\tau''\le T$ and reads the symbols
$x_{\kappa'}$ with $\kappa'\le\kappa''+1\le T+1$. The case $\kappa'=m$ would
require $m\le T+1$, which is excluded for $m\in\{n,n'\}$; hence the left
border is never read, and the input word is read only at positions
$\kappa'\le T+1$, where $J_n$ and $J_{n'}$ coincide. The two recursions are
therefore identical step by step, including the points where an argument is
$\Fq$ or undefined.
\end{proof}

The input word on $[0,T+1]$ consists of depth-row cells of depths at most
$2T+1$. These rows are driven only by the rows of smaller depth
\eqref{eq:comoving}, so the vector of the first $2T+2$ depth rows evolves
autonomously, and one period of it is a certificate for all later times.

\begin{lemma}[periodicity certificate]\label{lem:certificate}
Let $\delta$ be a rule, $D\ge0$, and $V_t:=(\Phi_t(0),\dots,\Phi_t(D))$ for
$t\ge D$. There is a map $G_D$ on $W^{D+1}$, depending on $\delta$ and $D$
only, with $V_{t+1}=G_D(V_t)$ whenever $t\ge D$ and $V_t\in W^{D+1}$. If
$V_{t_0}=V_{t_0+p}$ for some $t_0\ge D$ and $p\ge1$, and $V_t\in W^{D+1}$ for
$t_0\le t<t_0+p$, then $V_t\in W^{D+1}$ and $V_{t+p}=V_t$ for all $t\ge t_0$.
\end{lemma}

\begin{proof}
For $t\ge D$ and $0\le j\le D$ we have $j\le t$, so the arguments in
\eqref{eq:comoving} are $\Phi_t(j)$, $\Phi_t(j-1)$ and $\Phi_t(j-2)$, with
indices at most $D$, and $\Lq$ for negative indices; none of them is the border
symbol $\Phi_t(t+1)=\ast$. Hence $G_D(v)_j:=\delta(v_j,v_{j-1},v_{j-2})$, with
$v_{-1}=v_{-2}=\Lq$, has the required property. For the second statement,
suppose that $V_{t'}\in W^{D+1}$ and $V_{t'+p}=V_{t'}$ for $t_0\le t'\le t$,
which is true for $t=t_0$. Then $V_{t+1+p}=G_D(V_{t+p})=G_D(V_t)=V_{t+1}$, and
$V_{t+1}\in W^{D+1}$, by hypothesis if $t+1<t_0+p$ and because
$V_{t+1}=V_{t+1-p}$ otherwise. Induction on $t$ concludes.
\end{proof}

\begin{corollary}[finitely many early rows]\label{cor:reduce}
Let $\delta$, $T\ge0$, $D:=2T+1$, $t_0\ge D$ and $p\ge1$ satisfy the hypothesis
of \cref{lem:certificate}. Then $J_n(\kappa)=J_{n+p}(\kappa)\in W^2$ for all
$0\le\kappa\le T+1$ and $n\ge t_0+3$. Hence, for every $N_0\ge T+2$, every
$R_n$ with $n\ge N_0$ coincides on $E(T)$, in the sense of \cref{lem:early},
with $R_{n'}$ for one of the finitely many lengths $N_0\le
n'<\max(N_0,t_0+3)+p$.
\end{corollary}

\begin{proof}
For $1\le\kappa\le T+1$ the components of $J_n(\kappa)$ are cells of the depth
rows $2\kappa-2$ and $2\kappa-1\le D$ at the times $n-3+\kappa$ and
$n-2+\kappa$, which are at least $n-2>t_0$; by \cref{lem:certificate} they lie
in $W$ and do not change when $n$ is replaced by $n+p$. For $\kappa=0$ the
lower component is $\Lq$, by convention or by \cref{lem:lightcone}, and the
upper one is $C_\infty(n-2,n)=\Lq$. If $n\ge\max(N_0,t_0+3)+p$, subtracting a
multiple of $p$ gives a length $n'$ in $[\max(N_0,t_0+3),\max(N_0,t_0+3)+p)$
with $J_{n'}=J_n$ on $[0,T+1]$, and \cref{lem:early} applies since
$n,n'\ge T+2$.
\end{proof}

For a minimal-time solution, the depth rows are eventually periodic
(\cref{thm:third}), so $t_0$ and $p$ as in \cref{cor:reduce} always exist.
Thus the condition that no reflected triangle $R_n$ with $n\ge N_0$ has a
firing cell in its first $T+1$ rows reduces to finitely many lengths, with a
certificate $(t_0,p)$ found by simulating the half-line.

\subsection{Early-Row Formulas}

For a search over all rules, the half-line is not known in advance, but its
first rows and the lines of small lengths can be encoded together. The
following conditions collect what a minimal-time solution must satisfy on the
complete lines of small lengths, on the early rows and return chains of the
next lengths, and on the cone of their right ends.

\begin{definition}\label{def:G}
Let $k\ge3$, $T\ge1$, $N_1\ge T+2$ and $T+1\le C\le N_1$. A $k$-state rule
$\delta$ satisfying \eqref{eq:quiescence} \emph{satisfies $\G_k(T,N_1,C)$} if
\begin{enumerate}[(G1)]
\item $\delta$ synchronizes the line of length $n$ at time $2n-2$ for $2\le
n\le T+1$;
\item for $T+2\le n\le N_1$, $1\le i\le n$ and $t\le n-1+T$: $C_n(t,i)\in W$;
\item for $T+2\le n\le N_1$ and $1\le i\le n$ with $t:=2n-1-i\le n-1+T$:
$C_n(t,i)\ne C_\infty(t,i)$;
\item $C_\infty(t,i)\in W$ whenever $t+i\le 2N_1$;
\item for $T+2\le n\le C$: $R_n(\tau,\kappa)\in W$ for all $(\tau,\kappa)\in
R_n$ with $\tau+\kappa\le n-1$ and $\tau<n-1$, $R_n(n-1,0)=\Fq$, and
$R_n(\kappa,\kappa)\ne C_\infty(n-1+\kappa,n-\kappa)$ for $2\kappa\le n-1$.
\end{enumerate}
We write $\G_k(T,N_1)$ for $\G_k(T,N_1,T+1)$, in which (G5) is void.
\end{definition}

Condition (G1) asks for the complete lines of lengths up to $T+1$. Condition
(G2) says that no cell of the lines of lengths $T+2,\dots,N_1$ fires up to the
relative time $\tau=T$, and (G3) is the return chain of \cref{lem:return},
restricted to the relative times $\tau\le T$. Condition (G4) says that the
half-line does not fire on the cells with $t+i\le 2N_1$. Condition (G5)
concerns the backward dependency cone $\{(\tau,\kappa)\in R_n:\tau+\kappa\le
n-1\}$ of the firing cell $(n-1,0)$; by the local definition of $C_n$, its
cells are defined even if a cell outside the cone fires earlier.

\begin{lemma}\label{lem:G}
Every $k$-state minimal-time solution satisfies $\G_k(T,N_1,C)$ for all $T$,
$N_1$ and $C$ as in \cref{def:G}.
\end{lemma}

\begin{proof}
(G1) holds by definition. For (G2), $n\ge T+2$ gives $n-1+T<2n-2$, and every
cell of the line of length $n$ is in $W$ before time $2n-2$. (G3) is
\cref{lem:return}, which applies because the lines of lengths $n$ and $n+1$
are synchronized at times $2n-2$ and $2n$. (G4) is \eqref{eq:neverfire}. For
(G5), the cells of $R_n$ with $\tau<n-1$ are in $W$,
$R_n(n-1,0)=C_n(2n-2,n)=\Fq$, and the return chain differs from $C_\infty$ by
\cref{lem:return}.
\end{proof}

For a fixed half-line, the early rows of all lengths $n\ge N_0$ can thus be
checked on the representatives of \cref{cor:reduce}; \cref{prop:nearmiss} is
proved in this way, and the same representatives allow a formula to impose (G2)
and (G3) on all $n\ge N_0$.
